% !TeX root = ../../CursoFUSAL_Bloque0_revision.tex
% =====================================================================
%  Nota de revisión
%  SOLO aparece en el extracto que se manda a revisar (no en main.tex)
% =====================================================================
\chapter*{Cómo revisar este documento}
\addcontentsline{toc}{chapter}{Cómo revisar este documento}
\markboth{Cómo revisar este documento}{Cómo revisar este documento}

\section*{Qué es esto}

Este documento es un \textbf{extracto} del \textit{Curso de Ingeniería \FUSAL}:
contiene únicamente la parte I, el tronco común, que es la que se lee entera todo
el equipo antes de tocar el CAD. El libro completo tiene ocho partes y sesenta
capítulos; el resto todavía está sin redactar.

La numeración de capítulos y las referencias cruzadas son las del libro completo.
Por eso verás llamadas a capítulos que aquí no aparecen (por ejemplo, «capítulo
\ref{ch:evento_costes}»): son correctas, pero remiten a partes todavía no escritas.

\section*{Estado de cada capítulo}

\begin{center}
\begin{tabularx}{\textwidth}{@{}clX@{}}
\toprule
\textbf{Cap.} & \textbf{Título} & \textbf{Estado} \\
\midrule
\ref{ch:formula_student} & Formula Student: el juego y sus reglas & Redactado \\
\ref{ch:coche_como_sistema} & El coche como sistema & Redactado, breve a propósito \\
\ref{ch:proceso_diseno} & Proceso de diseño & \textbf{Solo los diagramas.} El texto está pendiente \\
\ref{ch:herramientas_cad_datos} & Herramientas: CAD, datos y versiones & Redactado \\
\ref{ch:mecanica_aplicada} & Fundamentos de mecánica aplicada & Redactado \\
\ref{ch:elementos_finitos} & Análisis por elementos finitos & Redactado \\
\ref{ch:ensayo_adquisicion_datos} & Ensayo y adquisición de datos & \textbf{Vacío a propósito}, se explica por qué \\
\bottomrule
\end{tabularx}
\end{center}

\section*{Sobre qué queremos que opines}

\begin{enumerate}
  \item \textbf{Errores de contenido.} Cualquier cosa que esté técnicamente mal,
        mal explicada o desactualizada. Es lo más importante.
  \item \textbf{Lo que no encaja con cómo trabajamos.} Buena parte de este
        material son \emph{propuestas}, no decisiones tomadas. En particular:
        \begin{itemize}
          \item la nomenclatura de piezas del apartado
                \ref{sec:herramientas_cad_datos:03} y la lista de códigos de
                subsistema;
          \item el sistema de coordenadas del apartado
                \ref{sec:herramientas_cad_datos:02};
          \item los pasos en \textcolor{fusalSec}{rojo discontinuo} de la figura
                \ref{fig:proceso_diseno}, que son añadidos al proceso de diseño y
                hay que decidir si se adoptan.
        \end{itemize}
        Si algo no te sirve, dilo ahora: cuesta mucho menos cambiarlo antes de
        que el equipo empiece a usarlo.
  \item \textbf{Lo que falta.} Temas que darías por imprescindibles en el tronco
        común y no están.
  \item \textbf{Si se entiende.} El público objetivo es alguien que acaba de
        entrar en el equipo. Si un apartado da por sabido algo que un rookie no
        sabe, señálalo.
\end{enumerate}

\section*{Sobre qué NO hace falta que opines}

\begin{itemize}
  \item \textbf{Los recuadros amarillos.} Son anotaciones de trabajo: cosas que ya
        sabemos que faltan. Están recogidas todas juntas al final del documento.
  \item \textbf{Los recuadros \textsc{caso fusal} que están sin rellenar.}
        Necesitan datos del equipo (puntuaciones, piezas rotas, masas) que
        todavía hay que recopilar.
  \item \textbf{Los capítulos \ref{ch:proceso_diseno} y
        \ref{ch:ensayo_adquisicion_datos}}, por lo dicho arriba. Del
        \ref{ch:proceso_diseno} sí interesa que revises las dos figuras.
  \item \textbf{La maquetación.} Portada, tipografía y colores están sin cerrar.
\end{itemize}

\section*{Cómo mandar los comentarios}

Indica siempre \textbf{número de capítulo y de apartado} (por ejemplo,
«\ref{ch:mecanica_aplicada}.4, lo del par de apriete») y qué cambiarías. Un
comentario sin localizar cuesta más de colocar que de escribir.

\vspace{0.8em}
\noindent
\begin{tabularx}{\textwidth}{@{}lX@{}}
\textbf{Enviar a:} & \pendiente{Rellenar: nombre y correo} \\[0.4em]
\textbf{Fecha límite:} & \pendiente{Rellenar} \\
\end{tabularx}
